\section{Introduction}
\label{sec:introduction}

3D Gaussian Splatting \citep{kerbl2023_3dgs} optimizes an explicit set of anisotropic 3D
Gaussians against photographs by differentiable rasterization.  Every iteration samples one
or more source views and compares a full render against the stored photograph, so all
images have to be available during the whole optimization.  Working with dense multi-camera
captures, this becomes the bottleneck.  The source images themselves are the largest
resident data structure of the reconstruction.  A camera dome producing tens to hundreds of
high-resolution views per frame quickly exceeds what a single GPU can hold, and streaming
the images from disk instead trades the memory bottleneck for a storage-bandwidth
bottleneck.  The model being optimized, a few hundred thousand Gaussians, is small by
comparison.

The optimization itself does not consume images.  It consumes colors at projected sample
locations.  Any representation that can answer which color view $v$ sees at image location
$\bm{x}$ can thus supervise the optimization.  We therefore fit every input image once,
offline, with a compact mixture of anisotropic 2D Gaussians, called a compact capture, and
discard the image.  The capture is three orders of magnitude smaller than the raw RGB
working set, see \cref{sec:memory}, analytic, differentiable and can be evaluated at
arbitrary subpixel locations without decoding, resampling or caching full frames.  Our
working hypothesis is that compact 2D Gaussian captures contain sufficient information to
supervise reconstruction and refinement of a standard 3D Gaussian Splatting representation
without returning to source RGB.

The resulting pipeline has a hard, checkable data boundary.  After the per-view fit, no
stage of reconstruction opens an image file, imports an image loader or receives a dense
RGB tensor.  Reconstruction consumes only calibrated cameras plus frozen 2D Gaussian
fields.  We treat this process separation as part of the result and enforce it
structurally through typed inputs, import denials and containment invariants, see
\cref{sec:overview}.\evidence{ara/logic/claims.md C31}

In this paper we therefore describe a supervision method that trains standard 3DGS from
frozen 2D Gaussian fields, together with a compact capture format with exact evaluation
semantics, see \cref{sec:fields,sec:reconstruction}.  The sampled objective is unbiased
and its sparse point renderer matches a dense reference renderer within verified
tolerances.  We define a controlled protocol which measures the memory of dense versus
compact supervision at matched initialization, seeds and stopping rules, see
\cref{sec:memory}.  \todo{Run the protocol and report the memory and quality trade-off.
This is the main missing evidence of the paper, see T-1 to T-6 in \cref{sec:todo-list}.}
Since rendering projects the 3D Gaussian field to per-view 2D Gaussians, the
reconstruction problem also admits a more direct initialization.  Inverting the projection
turns every fitted 2D Gaussian into a constraint on the 3D field.  We describe this
tomographic initialization, called Beam Fusion, in closed form, together with an
observability argument which fixes the minimum number of contributing views at three, see
\cref{sec:tomography}.  The naive inversion fails in identifiable ways.  We analyze these
failures and address them with a carrier refinement consisting of a renderer-aware
covariance repair, a two-phase fixed-topology schedule and a silhouette containment rule,
selected by paired ablations, see \cref{sec:refinement}.  \todo{Replicate the ablations on
fresh scenes with held-out cameras, see T-10.}

We focus on reconstruction from calibrated dense captures.  The reconstruction result is
established with the standard random and SfM initializations alone, so the value of the
tomographic initializer is a separate question with its own evaluation, and a weak
initializer outcome does not touch the reconstruction result.  \Cref{sec:experiments}
evaluates the full pipeline and closes with the ablation study.  Throughout, black text is
mechanism, definition or an evidence-bound result with the binding artifact referenced,
and red text is a TODO.  \Cref{sec:todo-list} collects all TODOs in one table.
